\cotefile{79}
\scanpage{205}
\mspage{105}

On l'obtient en relation avec la structure de graphe \inserted{$(S,A_e)$} sur $S$, associée au choix d'une base \inserted{de $\det M=\Lambda^2M$} ou ``jauge'' de $M$

\[
e\in\det(M)^* \tag*{(34)}
\]

par la \deleted{condition} formule

\[
A_e=\bigl\{\{s,t\}\in\mathscr P_2(S)\mid \inserted{\exists x,y\in M^*,\ \text{t.q.}}\ x\wedge y=e\ \deleted{\unsure{\dot x=s,\dot y=t}}\ \inserted{\dot x=s,\dot y=t}\bigr\} \tag*{(35)}
\]

(ne dépend que de la structure $\frac12$ symplectique de $M$ $\alpha=\{e,-e\}$ définie par $e$). Posant

\[
\left\{\begin{aligned}
&x\wedge y=\psi_e(x,y)\,e\\
&\psi_e\ \text{forme bilinéaire alternée non dég. définie par }e
\end{aligned}\right. \tag*{(36)}
\]

On voit donc que l'on a, pour $s\in S$, $s=\dot x$

\[
\operatorname{Et}_{A_e}(s)=
\{\deleted{\unsure{t\in S}}\inserted{\dot y}\mid\inserted{y\in M^*,}\ \psi_e(x,y)=\unsure{1|{\pm}1}\ \deleted{\unsure{\dot y=t}}\}. \tag*{(37)}
\]

Considérons alors l'application linéaire

\[
\left|\begin{aligned}
&\Phi_e:M\xrightarrow{\sim}M^\vee \qquad\text{caractérisée par}\\
&\langle x,\Phi_e(y)\rangle=\psi_e(x,y)\qquad\text{i.e.}\quad x\wedge y=\langle x,\Phi_e(y)\rangle\cdot e
\end{aligned}\right. \tag*{(38)}
\]

elle induit, par passage au quotient

\[
\deleted{\text{on a donc}}\quad \varphi_e:S\longrightarrow S^\vee,\qquad\text{et on aura, pour }s\in S \tag*{(39)}
\]

\[
\operatorname{Et}_{A_e}(s)=\{t\in S\mid\langle s,\varphi_e(t)\rangle\in R_{\unsure{1|A}}\},\quad\text{où} \tag*{(40)}
\]

\[
\left\{\begin{aligned}
&R_1\subset S\times S^\vee,\\
&\mathord{''}\ \{\deleted{\unsure{(s,t)}}\inserted{(\dot x,\dot x^\vee)}\}\in S\times S^\vee\mid \langle x,x^\vee\rangle=1\deleted{\unsure{t}}\}
=\{(s,s^\vee)\in S\times S^\vee\mid\inserted{\langle s,s^\vee\rangle=\dot 1}\}
\end{aligned}\right. \tag*{(41)}
\]

% slanted
\margin{où on introduit l'accouplement $(42)\quad S\times S^\vee\longrightarrow\Lambda/{\pm\mathrm{id}_\Lambda}$ par passage au quotient à partir de $M^*\times M^{\vee *}\longrightarrow\Lambda$.}

\deleted{On voit donc que les} \inserted{\unsure{Supposons|Supposant}} les structures de graphe $(S,A_e)$ \unsure{stricts|strictes}, i.e. qu'on a

\[
\left\{\begin{aligned}
&\forall(x,x')\in M^*\times M^*,\quad\text{la relation}\\
&\left[\begin{aligned}
&\forall y\in M^*\quad x\wedge y\in(\det M)^*\} \iff x'\wedge y\in(\det M)^*\\
&\hspace{38mm}\implies x'\wedge y\in\{\pm x\wedge y\}
\end{aligned}\right]\\
&\text{implique}\quad x'\in\pm x,\quad\deleted{\text{\unsure{i.e. que }} a\in\unsure{\mathrm{Fix}(M^*)},}\\
&\deleted{\exists\,\unsure{x,x'\in a}\ \text{et que }\unsure{\in}}
\end{aligned}\right. \tag*{(\unsure{43})}
\]

ce qui, pour $\Lambda$ $\frac12$ local, signifie simplement \unsure{qu'on|qu'un|qu'en}

% archive sheet number: 204
